\section{Part II --- Recognition-mode landings}\label{sec:recognition}

\subsection{Two columns}\label{sec:recognition:column-classification}

Part~I described how internal derivations stall.  Part~II describes the
other route used by the companion applications: a closure that names
its load-bearing source as a premise, instead of presenting it as the
output of an internal derivation.  We call such a closure a
\emph{\RML}.  This part is a vocabulary together with four named
patterns, stated as conditional schemas.  Their predicates---whether a
template applies, whether content is derivable from the framework's
primitives, whether a closure lands at theorem grade---are primitive, so
the schemas have no truth conditions until those predicates are given a
meaning; they describe an intended informal reading and are not
mathematical conjectures in the strict sense.  The two applications
that illustrate them are described in \Cref{sec:validation}.  Unlike
Part~I, the targets here are propositions rather than claims; only
\Cref{thm:attack-foreclosure-v5}, which refers back to Part~I, uses the
interpretation of a claim.

The organizing distinction comes from the V-Differential classification
of the hiddenness companion \cite{Tsiokos2026CSL}, which asks how a
substrate's witness is exposed to current observation.  In a
witness-content-exposing substrate, current observation can read a
witness datum, such as an assignment or a solution coordinate.  In a
trace-state-only substrate, current observation sees only trace
bookkeeping---step counts, audit fields and similar records of how a
computation proceeded---and no witness datum.

\begin{definition}[Columns]\label{def:vdiff-tso-column}
A substrate \(S\) lies in the \emph{trace-state-only column},
\(\operatorname{TSO}(S)\), when its load-bearing witness is not exposed to
current observation and is available only through trace-state evidence,
which a closure must name and accept as a premise.  It lies in the
\emph{witness-content-exposing column}, \(\operatorname{WCE}(S)\), when its
witness is exposed to current observation through the substrate's own
apparatus.  Both are primitive predicates.
\end{definition}

The hiddenness companion places the Navier--Stokes substrate in the
second column and the Riemann-hypothesis and P-versus-NP substrates in
the first.  It adopts these placements as hypotheses, with an informal
motivation for each, and this paper uses them with the same status.
For the Riemann hypothesis the named source is
\(\Gamma_{\mathrm{SDTC}}^{\mathrm{Selberg}}\), recorded in the
application \cite{Tsiokos2026RH}; the self-dual trace confinement
companion \cite{Tsiokos2026SDTC} supplies the squeeze it feeds.  Its name refers to Selberg's trace
formula \cite{Selberg1956TraceFormula} by analogy only: Selberg's
formula concerns automorphic spectra and does not provide the
completed-zeta data that the source asserts, which remain a premise.
For \(P\) versus \(NP\) the named source is
\(\Gamma_{\mathrm{csl\text{-}sat\text{-}hidden}}\), recorded in the
application \cite{Tsiokos2026PvNP}.

Naming a source is not a sufficient premise in itself.  The
Riemann-hypothesis endpoint also needs an identification of its
coordinates with the classical zeros, including the conditions that
the identification preserves the real part \(\frac12\) and covers every
zero; the Navier--Stokes endpoint needs its spectral closure and its
hypotheses on the initial data; and the P-versus-NP endpoint needs an
accepted negative source and, in its family form, the admission of
search events.  Recognition-mode landings and derivation-mode landings
are two different kinds of conditional closure; neither is a weaker
form of the other.

\subsection{The seven-stage template}\label{sec:recognition:template-applicability}

The template describes a recognition-mode landing as a sequence of
seven checks on a formed closure \(c\) with target \(X\).
\begin{enumerate}
\item \emph{Construction.}  Build the formed closure with its types,
  levels, thresholds and visibility data, and exhibit that it reaches
  beyond the carriers available before it was formed.
\item \emph{Translation.}  Prove that the closure's internal predicate
  is equivalent to the standard statement of \(X\).  This is not where
  the source enters; it ensures that an accepted source yields the
  standard statement.
\item \emph{Applicability.}  Check that the relevant structural theorem
  applies to the closure, and separate the hypotheses already verified
  from those still needed.
\item \emph{Smuggling check.}  Check that the admissibility,
  lawfulness and saturation hypotheses do not already assume the source.
\item \emph{Anti-tautology.}  Check that the closure's predicate
  genuinely depends on the closure's data, so that the apparent gain
  does not split into independent claims about independent carriers.
\item \emph{Derivation sweep.}  Search for a derivation of the source
  along three routes (\Cref{def:three-option-sweep}).
\item \emph{Recognition.}  Name the source and record its acceptance as
  a separate premise.
\end{enumerate}
We write \(T=\langle S_{1},\ldots,S_{7}\rangle:\mathsf{7Stage}(c,X)\) for a
record of the seven stages as propositions about \((c,X)\).

\begin{definition}[Template applies]\label{def:applies}
For a formed closure \(c\), a target \(X\) and a seven-stage record
\(T\), \(\operatorname{Applies}(c,X,T)\) is a primitive predicate,
intended to state that the seven checks of \(T\) succeed for \((c,X)\);
it is not defined as the conjunction of the seven stages.
\end{definition}

\begin{schema}[Recognition-Mode Template Applicability]\label{thm:template-applicability}
Let \(c\) be a formed closure with target \(X\), and suppose that
\(\operatorname{TSO}(c)\), that the SB closure assumption holds for \(c\),
and that \(T:\mathsf{7Stage}(c,X)\) satisfies
\(\operatorname{Applies}(c,X,T)\).  Then \(c\) lands at theorem grade:
\(\operatorname{TheoremGrade}(c,X)\).
\end{schema}

All four inputs are needed; in particular the schema assumes, and does
not derive, that the template applies.  \(\operatorname{TheoremGrade}\)
is primitive, so even an instance of the schema would not prove \(X\);
it would record that the closure has the status of a conditional
theorem under its premises.  The two applications of
\Cref{sec:validation} illustrate the schema.

\subsection{The derivation sweep}\label{sec:recognition:verdict-signature}

\begin{definition}[Three-option derivation sweep]\label{def:three-option-sweep}
The \emph{derivation sweep} of Stage~6 tests three routes to the
load-bearing source: the \emph{framework-primitive option}, a
derivation from the primitives of the framework; the \emph{non-descent
option}, whether a derivation is blocked because the needed content
does not descend to the coarser layer; and the \emph{elevation option},
whether passing to the V-Differential classification resolves the
route.  Each option receives one of three verdicts: (i) the option is
the closure route; (ii) the option is not the closure route but does
not block it; (iii) the option passes the question to the
V-Differential classification.
\end{definition}

Formally, \(\mathsf{SweepOption}=\{\mathsf{FP},\mathsf{ND},\mathsf{VDE}\}\),
\(\mathsf{Verdict}=\{\mathrm{i},\mathrm{ii},\mathrm{iii}\}\), and a
seven-stage record \(T\) has a primitive verdict map
\(v_{T}:\mathsf{SweepOption}\to\mathsf{Verdict}\).

\begin{schema}[Cross-Track Verdict Signature]\label{thm:verdict-signature}
Let \(c\) be a formed closure in the trace-state-only column with
target \(X\), and let \(T:\mathsf{7Stage}(c,X)\) satisfy
\(\operatorname{Applies}(c,X,T)\).  Then
\[
  v_T(\mathsf{FP})=\mathrm{ii},\qquad
  v_T(\mathsf{ND})=\mathrm{iii},\qquad
  v_T(\mathsf{VDE})=\mathrm{ii}.
\]
\end{schema}

The intended explanation is that a verdict (i) for the
framework-primitive option would contradict
\Cref{thm:closure-content-thesis}, and that the other two options only
relabel the route; this is a heuristic, not a proof.

The evidence is limited to the two trace-state-only applications.  The
first version of the Riemann-hypothesis application reports that three
candidate derivations of its source from the framework's primitives
each fail \cite{Tsiokos2026RHvOne}, and the first version of the
P-versus-NP application reports the same for its source
\cite{Tsiokos2026PvNPvOne}.
\Cref{tab:verdict-signature} records this paper's reading of those
reports as verdicts.  The assignment is not computed by any formal
procedure, and no formal instance of the schema is constructed.

\begin{table}[H]
  \centering
  \caption{Verdicts of the derivation sweep, as read in this paper from
  the reports of the two applications; not a formal computation.}
  \label{tab:verdict-signature}
  \small
  \begin{tabularx}{\textwidth}{@{}>{\raggedright\arraybackslash}p{0.22\textwidth}>{\centering\arraybackslash}X>{\centering\arraybackslash}X>{\centering\arraybackslash}X@{}}
    \toprule
    Application & Framework-primitive option & Non-descent option & Elevation option \\
    \midrule
    \(P\) versus \(NP\) & (ii) & (iii) & (ii) \\
    Riemann hypothesis & (ii) & (iii) & (ii) \\
    \bottomrule
  \end{tabularx}
\end{table}


\subsection{The closure-content thesis}\label{sec:recognition:closure-content-thesis}

The derivation sweep suggests a semantic distinction.  On a
trace-state-only substrate the load-bearing content is part of the
formed closure itself: the closure names it and accepts it as a
premise.  We use three primitive predicates of a structural fact
\(s\): \(\operatorname{FLC}(s,c)\), that \(s\) is \emph{formed-layer
content} of \(c\), supplied by the formed closure rather than produced
later; \(\operatorname{DFP}(s)\), that \(s\) is \emph{derivable from
the framework's primitives}; and \(\operatorname{NAS}(s,c)\), that \(s\)
is \emph{named as an accepted source} of \(c\).

\begin{definition}[Load-bearing content]\label{def:load-bearing}
For a formed closure \(c\) and a target \(X\), the load-bearing content
\(s_{c,X}=\operatorname{LoadBearingFor}(c,X)\) is the structural fact
that \(c\) must accept as a premise in order to land \(X\).  The
assignment \((c,X)\mapsto s_{c,X}\) is primitive.
\end{definition}

\begin{schema}[Closure-Content Thesis]\label{thm:closure-content-thesis}
Let \(c\) be a formed closure in the trace-state-only column, with
target \(X\).  Then its load-bearing content is formed-layer content,
is not derivable from the framework's primitives, and is named as an
accepted source:
\[
  \operatorname{FLC}(s_{c,X},c)\wedge
  \neg\operatorname{DFP}(s_{c,X})\wedge
  \operatorname{NAS}(s_{c,X},c).
\]
\end{schema}

The three predicates are independent: a fact can be formed-layer
content without being named, be named without being load-bearing, or be
derivable without being a recognition source.  The thesis asserts the
specific combination for the load-bearing content of a trace-state-only
closure.  It concerns the content \(s_{c,X}\) only; that every
recognition-mode landing of \(c\) imports this content is a separate
clause of \Cref{thm:attack-foreclosure-v5}.

\subsection{Refined attack foreclosure}\label{sec:recognition:attack-foreclosure-v5}

The last schema combines the two parts.  Write
\(\operatorname{DerivationRoute}(m)\) for the primitive predicate that
an internal move \(m\) is a derivation route,
\(l:\operatorname{RMLanding}(c,X)\) for a recognition-mode landing of
\(X\) on \(c\), and \(\operatorname{ImportsAsNamedRecognition}(l,s)\)
for the primitive predicate that \(l\) imports \(s\) as a named source.

\begin{schema}[Refined Attack Foreclosure]\label{thm:attack-foreclosure-v5}
Let \(x\) be a claim in framework scope, put \(X=\llbracket x\rrbracket\),
and let \(c\) be a formed closure in the trace-state-only column for which
the SB closure assumption holds.  Then
\[
\begin{aligned}
  &\forall m.\;
    \operatorname{DerivationRoute}(m)
    \Longrightarrow
        \neg\operatorname{AdvancesWithoutExternal}(m,x),\\
  &\forall l:\operatorname{RMLanding}(c,X).\;
    \operatorname{ImportsAsNamedRecognition}(l,s_{c,X}),\\
  &\forall l:\operatorname{RMLanding}(c,X).\;
    \operatorname{TheoremGrade}(c,X).
\end{aligned}
\]
\end{schema}

The first clause is a consequence of \Cref{thm:attack-foreclosure-v4},
restricted to derivation routes.  The second and third clauses are new:
a recognition-mode landing imports the load-bearing content as a named
source, and lands at theorem grade.  The schema therefore adds to
\Cref{thm:attack-foreclosure-v4} and does not retract it.  It does not
assert that every admissible route is either a derivation route or a
recognition-mode landing; that the two kinds exhaust the routes is an
additional assumption, which we do not make.

\begin{table}[H]
  \centering
  \caption{The four statements of Part~II.  All are conditional schemas
  with primitive predicates; none is proved.}
  \label{tab:recognition:summary}
  \small
  \renewcommand{\arraystretch}{1.1}
  \begin{tabularx}{\textwidth}{@{}>{\raggedright\arraybackslash}p{0.27\textwidth}>{\raggedright\arraybackslash}X@{}}
    \toprule
    Statement & Content \\
    \midrule
    \Cref{thm:template-applicability} & On a trace-state-only formed closure satisfying the SB closure assumption, a seven-stage record for which the template applies yields theorem-grade landing. \\
    \Cref{thm:verdict-signature} & Under the same hypotheses, the derivation sweep returns the verdicts \SigPattern. \\
    \Cref{thm:closure-content-thesis} & The load-bearing content of a trace-state-only closure is formed-layer content, is not derivable from the primitives, and is named as an accepted source. \\
    \Cref{thm:attack-foreclosure-v5} & For an in-scope target and a trace-state-only formed closure satisfying the SB closure assumption: derivation routes do not advance the target without outside content, and recognition-mode landings import the load-bearing content and land at theorem grade. \\
    \bottomrule
  \end{tabularx}
\end{table}

